\sheettitle
  {Independent solutions}
  {Expectation, laws, tail integrals, and finite conditioning}

Use this document only after a complete attempt. Compare the first step at which your argument ceases to be legitimate; do not replace the work record with a cleaner proof.

\section*{Exercise 0 -- The $L^p$--RN retention gate}

Since $\mathbb P(\Omega)=1$ and $q>1$, the finite-measure embedding gives
\[
\|Y\|_1\le\|Y\|_q,
\]
so $Y\in L^1(\mathbb P)$. The numerator is the restriction of the finite signed measure $Y\mathbb P$ to $(\Omega,\mathcal G)$; the reference is $\mathbb P|_{\mathcal G}$; both live on $(\Omega,\mathcal G)$. Moreover,
\[
|\nu_Y^{\mathcal G}|(\Omega)
\le
\mathbb E[|Y|]
<+\infty,
\]
and $\mathbb P(A)=0$ implies $\nu_Y^{\mathcal G}(A)=0$.

The signed Radon--Nikodym theorem therefore supplies a $\mathcal G$-measurable
\[
H\in L^1(\mathbb P|_{\mathcal G})
\]
such that
\[
\nu_Y^{\mathcal G}(A)
=
\int_AH\,d\mathbb P,
\qquad A\in\mathcal G,
\]
with $H$ unique $\mathbb P$-almost surely. This representative is $\mathbb E[Y\mid\mathcal G]$.

For a random element $X$, LOTUS is the pushforward integration identity for the measure $\mu_X=X_\#\mathbb P$. It therefore needs no dominating Lebesgue measure. A formula involving $f_X$ additionally requires $\mu_X\ll\lambda$; only then does Radon--Nikodym provide $f_X=d\mu_X/d\lambda$, unique $\lambda$-almost everywhere.

\section*{Exercise 1 -- Expectation through the law}

For $C\in\mathcal S$,
\[
\mathbb E[\ind_C(X)]
=
\mathbb P(X\in C)
=
\mu_X(C)
=
\int_S\ind_C\,d\mu_X.
\]
If $\varphi=\sum_{k=1}^na_k\ind_{C_k}$ is nonnegative simple, linearity gives
\[
\mathbb E[\varphi(X)]
=
\sum_{k=1}^na_k\mu_X(C_k)
=
\int_S\varphi\,d\mu_X.
\]
For nonnegative measurable $\varphi$, choose $\varphi_n\uparrow\varphi$ with $\varphi_n$ simple. Monotone convergence on $S$ and on $\Omega$ yields
\[
\mathbb E[\varphi(X)]
=
\lim_n\mathbb E[\varphi_n(X)]
=
\lim_n\int_S\varphi_n\,d\mu_X
=
\int_S\varphi\,d\mu_X.
\]
For signed $\varphi$, apply this result to $\varphi^+$ and $\varphi^-$ whenever their difference is defined.

On $((0,1),\mathcal B((0,1)),\lambda)$, the nonnegative variable $X(t)=1/t$ has infinite expectation. For an undefined signed expectation, take the identity variable on $\mathbb R$ under the standard Cauchy probability law: both $\mathbb E[X^+]$ and $\mathbb E[X^-]$ are infinite.

If $\mu_X$ is atomic, integration against it is a weighted sum. If $\mu_X\ll\lambda$ with density $f_X$, the same integral is represented by $\int\varphi(x)f_X(x)\,dx$. Neither representation changes the definition of expectation.

\section*{Exercise 2 -- Subgraphs, tail integration, and insurance payments}

For $q\in\mathbb Q_{>0}$, both $\{f>q\}\in\mathcal A$ and $(0,q)\in\mathcal B((0,+\infty))$. The identity
\[
H_f
=
\bigcup_{q\in\mathbb Q_{>0}}
\{f>q\}\times(0,q)
\]
therefore proves product measurability. Indeed, if $t<f(x)$, density of the rationals supplies $q$ with $t<q<f(x)$; the converse is immediate.

Because $\ind_{H_f}\ge0$, Tonelli applies without any prior integrability claim:
\[
\begin{aligned}
\int_Ef(x)\,\mu(dx)
&=
\int_E\int_0^\infty\ind_{\{t<f(x)\}}\,dt\,\mu(dx)\\
&=
\int_0^\infty\int_E\ind_{\{f(x)>t\}}\,\mu(dx)\,dt\\
&=
\int_0^\infty\mu(f>t)\,dt.
\end{aligned}
\]

Apply the same identity to the transformed variables. Since
\[
\mathbb P((X-d)_+>t)=\mathbb P(X>d+t),
\]
the change of variable $s=d+t$ gives
\[
\mathbb E[(X-d)_+]
=
\int_d^\infty\mathbb P(X>s)\,ds.
\]
Since $\mathbb P(X\wedge u>t)=\mathbb P(X>t)$ for $0\le t<u$ and is zero for $t\ge u$,
\[
\mathbb E[X\wedge u]
=
\int_0^u\mathbb P(X>t)\,dt.
\]
If $X\sim\operatorname{Exp}(\lambda)$, these are
\[
\frac{e^{-\lambda d}}{\lambda}
\quad\text{and}\quad
\frac{1-e^{-\lambda u}}{\lambda}.
\]

For the limited-loss payment,
\[
\mathbb E[\min\{(X-d)_+,u\}]
=
\int_d^{d+u}\mathbb P(X>s)\,ds.
\]
Under the exponential law this equals
\[
\frac{e^{-\lambda d}-e^{-\lambda(d+u)}}{\lambda}
=
\frac{e^{-\lambda d}(1-e^{-\lambda u})}{\lambda}.
\]

If $g$ is convex and both $X$ and $g(X)$ are integrable, Jensen gives
\[
g(\mathbb E[X])\le\mathbb E[g(X)].
\]
It supplies a bound. It does not evaluate the tail integral defining the expected payment.

\section*{Exercise 3 -- Conditioning on a finite partition}

Let $\mathcal G=\sigma(B_1,\ldots,B_r)$. A finite real-valued $\mathcal G$-measurable random variable is constant on each atom, so it has the form
\[
H=\sum_{i=1}^rc_i\ind_{B_i}.
\]
Testing the conditional-expectation identity on $B_j$ gives
\[
c_j\mathbb P(B_j)
=
\mathbb E[Y\ind_{B_j}].
\]
When $\mathbb P(B_j)>0$, this determines
\[
c_j
=
\frac{\mathbb E[Y\ind_{B_j}]}{\mathbb P(B_j)}.
\]
Hence
\[
\mathbb E[Y\mid\mathcal G]
=
\sum_{i=1}^r
\frac{\mathbb E[Y\ind_{B_i}]}{\mathbb P(B_i)}\ind_{B_i}.
\]
Taking expectations yields
\[
\mathbb E[Y]
=
\sum_{i=1}^r\mathbb P(B_i)\mathbb E[Y\mid B_i].
\]
On a null atom, the equation is $0=0$, so its coefficient is arbitrary. This is precisely version freedom on a null set. Taking $Y=\ind_A$ gives
\[
\mathbb P(A)
=
\sum_i\mathbb P(A\mid B_i)\mathbb P(B_i),
\]
the law of total probability.

\section*{Exercise 4 -- Bayes and representation bookkeeping}

The joint masses satisfy
\[
\mathbb P(A\cap B_i)
=
\mathbb P(A\mid B_i)\mathbb P(B_i).
\]
Summing over the partition gives
\[
\mathbb P(A)
=
\sum_i\mathbb P(A\mid B_i)\mathbb P(B_i).
\]
Therefore
\[
\mathbb P(B_j\mid A)
=
\frac{\mathbb P(A\mid B_j)\mathbb P(B_j)}
{\sum_i\mathbb P(A\mid B_i)\mathbb P(B_i)}.
\]

For $C\in\mathcal F$,
\[
\mathbb P_A(C)
=
\frac{\mathbb P(C\cap A)}{\mathbb P(A)}
=
\int_C\frac{\ind_A}{\mathbb P(A)}\,d\mathbb P.
\]
Thus
\[
\frac{d\mathbb P_A}{d\mathbb P}
=
\frac{\ind_A}{\mathbb P(A)}
\quad\mathbb P\text{-almost surely}.
\]

The two requested ledger rows are:

{\footnotesize
\hbadness=10000
\begin{tabularx}{\textwidth}{@{}Y Y Y Y Y@{}}
\toprule
Representative & Numerator measure & Reference & Domain & Uniqueness \\
\midrule
$d\mathbb P_A/d\mathbb P$ & $C\mapsto\mathbb P(C\cap A)/\mathbb P(A)$ & $\mathbb P$ & $(\Omega,\mathcal F)$ & $\mathbb P$-a.s. \\
$\mathbb E[Y\mid\mathcal G]$ & $C\mapsto\mathbb E[Y\ind_C]$, $C\in\mathcal G$ & $\mathbb P|_{\mathcal G}$ & $(\Omega,\mathcal G)$ & $\mathbb P$-a.s. \\
\bottomrule
\end{tabularx}
\par
}

If $X$ has a continuous law, $\mathbb P(X=x)=0$ for each fixed $x$. The event-ratio denominator therefore vanishes. A conditional kernel or density may provide a version, but it is a different construction and is not determined pointwise everywhere by the joint law.

\section*{Exercise 5 -- Original deployment analogues}

\textbf{(a) Capped-payment mechanism.}
The transformed payment is $X\wedge1$. The tail formula gives
\[
\mathbb E[X\wedge1]
=
\int_0^1\mathbb P(X>t)\,dt
=
\int_0^1(1+t)^{-2}\,dt
=
\left[-\frac1{1+t}\right]_0^1
=
\frac12.
\]

\textbf{(b) Nonlinear Poisson moment.}
For $N\sim\operatorname{Poisson}(2)$,
\[
\mathbb E[N]=2,
\qquad
\operatorname{Var}(N)=2,
\qquad
\mathbb E[N^2]=2+2^2=6.
\]
Hence
\[
\mathbb E\!\left[\frac{N(N+1)}2\right]
=
\frac{\mathbb E[N^2]+\mathbb E[N]}2
=
\frac{6+2}{2}
=
4.
\]
The required object is $\mathbb E[g(N)]$; $g(\mathbb E[N])=3$ is a different quantity.

\textbf{(c) Finite-mixture inversion.}
The joint masses are
\[
\mathbb P(A\cap B_1)=0.3\cdot0.8=0.24,
\qquad
\mathbb P(A\cap B_2)=0.7\cdot0.2=0.14.
\]
Thus $\mathbb P(A)=0.38$ and
\[
\mathbb P(B_1\mid A)
=
\frac{0.24}{0.38}
=
\frac{12}{19}.
\]

\section*{Closed-book reconstruction}

The first chain is
\[
X
\longmapsto
\mu_X=X_\#\mathbb P
\longmapsto
\mathbb E[\varphi(X)]=\int_S\varphi\,d\mu_X.
\]
No density is needed. A pdf appears only after choosing a reference $\lambda$ and verifying $\mu_X\ll\lambda$.

For a finite partition, $\mathcal G$-measurability forces atomwise constants; the identities
\[
\int_{B_i}H\,d\mathbb P
=
\int_{B_i}Y\,d\mathbb P
\]
determine those constants on positive-probability atoms. Summation yields total expectation, and normalization of the corresponding joint masses yields Bayes' formula. Null atoms retain arbitrary coefficients because conditional expectation is defined only up to almost-sure equality.
